\subsection{Vector Fields}
\lecdate{(check) Lec 13 - Feb 24 (Week 7)}
we jump arnd a bit - doing chs 6/9 kinda together.

whats a vector field? morally, a smooth choice of tangent vec at every pt in $M$.
eg, for $M\subseteq\bR^{n}$, tangent vecs live in $T_{p}M\subseteq\bR^{n}$.

\begin{defn}
	a \textbf{vector field} is a smooth map $X:M\gto\bR^{n}$ s.t.
	$X(p)\in T_{p}M\subseteq\bR^{n}$.
\end{defn}
what does it mean on an abstract mfld? ``lets start by cheating".
v1: look at charts.

\begin{defn}
	a \textbf{vector field} $X$ on $M$ is a map $p\mto X(p)\in T_{p}M$, such that
	for each chart $\vphi$, we have $D_{p}\vphi\circ X:U\gto\bR^{n}$ is smooth.
\end{defn}
this uses the identification $T_{p}M\simeq\bR^{n}$ induced by chart.
but what is that thing between $D_{p}\vphi\circ X$?

\begin{defn}
	the \textbf{tangent bundle} $TM$ of $M^{m}$ is a smooth $2m$-mfld with:
	\begin{itemize}
		\item pt set $\bigsqcup_{p\in M}T_{p}M$
		\item (non-max'l) atlas w charts of the form:
			\begin{equation*}
				\left(\bigsqcup_{p\in U}T_{p}M,(\vphi\circ\pi,D\vphi)\right) \qquad
				(\vphi\circ\pi,D_{p}\vphi):\bigsqcup_{p\in U}T_{p}M\gto\bR^{m}\times\bR^{m}
			\end{equation*}
	\end{itemize}
	for charts $(U,\vphi)$ in $M$.
\end{defn}
check: $\pi$ is smooth. remarks:
\begin{enumerate}[(1)]
	\item the tangent bundle comes w a proj $\pi:TM\gto M$, given by $T_{p}M\mto p$.
	\item for each $p\in M$, the preimg (or fiber) $\pi^{-1}(p)$ is a vspace.
	\item \textit{local triviality:} for each $p\in M$, there's a nbhd $U\ni p$ s.t. the diagram commutes:
		[TODO: add diagram], where $(D\vphi,\pi)$ is a smooth map and
		linear on each $\pi^{-1}(q)$.
\end{enumerate}

\lecdate{Lec 15 - Mar 5 (Week 8)}
accidentally git yeeted tuesdays notes so those are gone forever

last time: a VF $X$ on $M$ defns a flow $\Phi:J\subseteq\bR\times M\gto M$ where $\dot{\Phi}_{t}(p)=X_{\Phi_{t}(p)}$,
for some $J\ni\set{0}\times M$ open.

\begin{defn}
	a VF is \textbf{complete} if $J=\bR\times M$.
\end{defn}
for instance, $x\p_{x}$ is complete w $\Phi_{t}(x)=e^{tx}$, but
$x^{2}\p_{x}$ is not (escapes to $\infty$ in finite time).

\begin{thm}
	if $X$ has cpt sppt, it is complete.
\end{thm} \

\begin{pf} \vspace{-0.2in}
	\begin{enumerate}[1)]
		\item if $p\notin\supp(X)$, $X_{p}=0$ implies $\Phi_{t}(p)=p$ for all $t$.
			if $p\in\supp(X)$, then $\Phi_{t}(p)\in\supp(X)$ when defnd.
		\item since $J$ open, $\fa\ep>0$, there is an open $(-\ep,\ep)\times U_{\ep}$, where:
			\begin{equation*}
				U_{\ep} \ = \ \set{p\in M:\Phi_{t}(p) \trm{ defnd for } t \in (-\ep,\ep)}
			\end{equation*}
			these cvr $M$. by cptness, take finite subcvr $\set{U_{\ep_{i}}}$.
			then $U_{\ep_{\min}}$=M.
		\item by flow property, $\Phi_{\ep_{\min}/2}$ defnd $\fa p\in M$.
			thus $\Phi_{n\ep_{\min}/2}(p)$ is defined on $t\geq0$,
			and $h<0$ can be handled the same way.
	\end{enumerate}
\end{pf}
summary: if $X$ a complete VF, its flow defns a 1-param. family of diffeos.

\begin{crll}
	any VF on a cpt mfld is complete.
\end{crll}
is the crll out of order? oh well. note:
\begin{enumerate}[i)]
	\item $\Phi_{t}\circ\Phi_{s}=\Phi_{t+s}$
	\item $\Phi(t,p)$ is a smooth fn on $\bR\times M$
\end{enumerate}
so $\Phi(t,p)$ a smooth fn on $\bR\times M$.

\begin{prop}
	if $\Phi_{t}$ satisfies i) and ii) above, its the flow of a VF $X$, where:
	\begin{equation*}
		X_{p}(f) \ = \ \frac{\d}{\d t}\bigg\rvert_{t=0}f(\Phi_{t}(p))
	\end{equation*}
\end{prop}
in other words, $X_{p}$ is the vel vec of $\gamma(t)=\Phi_{t}(p)$.

\begin{pf}
	first, $X$ satisfies the product rule:
	\begin{equation*}
		X_{p}(fg) \ = \ \frac{\d}{\d t}\bigg\rvert_{t=0}f(\Phi_{t}(p))g(\Phi_{t}(p))
	\end{equation*}
	its a VF. trust. it is also an integral curve of $X$, as:
	\begin{equation*}
		\dot{\Phi}_{t}(p) \ = \ \frac{\d}{\d s}\bigg\rvert_{s=0}\Phi_{s}(\Phi_{t}(p)) \ = \ X_{\Phi_{t}(p)}
	\end{equation*}
	sure.
\end{pf}

example: given $A\in M_{n}(\bR)$, let $\Phi_{t}:\bR^{m}\gto\bR^{m}$ be:
\begin{equation*}
	\Phi_{t}(x) \ = \ e^{tA}x \ = \ \left(\sum_{n=0}^{\infty}\frac{t^{n}}{n!}A^{n}\right)x
\end{equation*}
then the corresp. VF is:
\begin{equation*}
	X(f)(x) \ = \ \frac{\d}{\d t}\bigg\rvert_{t=0}f(e^{tA}x)
	\ = \ \nabla f\cdot\frac{\d}{\d t}\bigg\rvert_{t=0}e^{tA}x
	\ = \ \nabla f\cdot Ax
\end{equation*}
so in ``wishy-washy physics notation", we have $X=Ax\cdot\nabla$.

now, an alt defn of the lie bracket.
we start with, naturally, the pullback.

\begin{defn}
	given $F:M\gto N$, we denote by $F^{*}f=f\circ F$ the \textbf{pullback} of $f:N\gto\bR$.

	if $F$ is a diffeo, then for a VF $X$ on $N$, we can define $F^{*}X$ by:
	\begin{equation*}
		F^{*}X_{F^{-1}(p)}=TF^{-1}(X_{p})
	\end{equation*}
\end{defn} \

\begin{defn}
	the \textbf{lie derivative} of $f:M\gto\bR$ wrt a VF $X$ is given as:
	\begin{equation*}
		L_{X}f \ = \ \frac{\d}{\d t}\bigg\rvert_{t=0}\Phi_{t}^{*}f \ \in \ C^{\infty}(M)
	\end{equation*}
	the lie derivative of a VF $Y$ wrt $X$ is:
	\begin{equation*}
		L_{X}Y \ = \ \frac{\d}{\d t}\bigg\rvert_{t=0}\Phi_{t}^{*}Y \quad \in \frk{X}(M)
	\end{equation*}
	where $\Phi$ is the flow of $X$.
\end{defn}
check: $L_{X}f=X(f)$.

\begin{thm}
	$L_{X}Y \ = \ [X,Y]$.
\end{thm}

\lecdate{Lec 16 - Mar 10 (Week 9)}
recall the defn of the lie deriv. of a VF.
since $\Phi$ is a diffeo, it induces an isomorphism
$T_{p}M\xrightarrow{ \ T_{p}(\Phi_{t}) \ }T_{\Phi_{t}(p)M}$.
this gives a way to compare vectors in these spaces, as $\Phi_{t}^{*}Y$ pulls
tangent vectors back into $T_{p}M$ (i.e. all vecs $\Phi_{t}^{*}Y\in T_{p}M$).

\begin{pf}[]
	[of 8.10] given $f\in C^{\infty}(M)$, wts $L_{X}Y(f)=[X,Y](f)$. we have:
	\begin{equation*}
		\Phi_{t}^{*}(Y(f)) \ = \ (\Phi_{t}^{*}Y)(\Phi_{t}^{*}f)
	\end{equation*}
	taking deriv at $t=0$, we get:
	\begin{equation*}
		X(Y(f)) \ = \
		\left(\frac{\d}{\d t}\bigg\rvert_{t=0}\Phi_{t}^{*}(Y)\right)(f)
		+ Y(\Phi_{t}^{*}f)
		\ = \ L_{X}Y(f) + Y(X(f))
	\end{equation*}
	thus $L_{X}Y=XY-YX=[X,Y]$ as needed.
\end{pf}
remark: $L_{X}f:=\displaystyle\frac{\d}{\d t}\bigg\rvert_{t=0}\Phi_{t}^{*}f$
measures the change in $f$ along an integral curve(?) of $X$, that is,
a curve whose velocity vector at $t=0$ is $X_{p}$.
this is precisely $X(f)$.
thus, $[X,Y]$ is how $Y$ changes along the flow of $X$.
in particular, we have the following:

\begin{thm}
	sps $\Phi$ flow of $X$, $\Psi$ flow of $Y$. then:
	\begin{equation*}
		[X,Y]=0 \ \iff \ \Phi_{t}^{*}Y = Y \ \iff \ \Psi_{t}^{*}X = X
		\ \iff \ \Phi_{t}\circ\Psi_{s} = \Psi_{s}\circ\Phi_{t}
	\end{equation*}
	for all $s,t$.
\end{thm} \

\begin{pf}
	the first two equivalences are clear from the definition.
	if $\Phi_{t}^{*}(Y)=Y$, then $\Phi_{t}$ takes integral curves of $Y$ to
	integral curves of $Y$. uhm.
	
	he literally just gave up lmao.
\end{pf}
more generally, we have that:
\begin{equation*}
	[X,Y] \ = \ \lim_{t\sto0}\frac{1}{t^{2}}[\Phi_{t}^{*}\Psi_{t}^{*}-\Psi_{t}^{*}\Phi_{t}^{*}]
\end{equation*}
lets not attempt to prove it this time.

and now, an interpretation of the jacobi identity. recall:
\begin{equation*}
	[X,[Y,Z]] + [Y,[Z,X]] + [Z,[X,Y]] \ = \ 0
\end{equation*}
we can rewrite this into the following form (which saany uses):
\begin{equation*}
	[X,[Y,Z]] \ = \ [[X,Y],Z] + [Y,[X,Z]]
\end{equation*}
this can be interpreted as a product rule, where $[\cdot,\cdot]$ is both
the derivation and the product.
``the lie bracket on derivations is a derivation on lie brackets".

a theorem! when can you take lin. indep. VF $X_{1},\dots,X_{r}$ and ``integrate them" to
get an $r$-dim submfld $S$?
i.e., for every $p\in M$, there is locally a unique $r$-dim submfld $S$ thru $p$
s.t. $X_{1},\dots,X_{r}$ are tangent to $S$ near $p$.

\begin{thm}[title=Frobenius' Theorem]
	the above happens iff:
	\begin{equation*}
		[X_{i},X_{j}] \ = \ \sum_{k=1}^{r}c_{ij}^{k}(p)X_{k}
	\end{equation*}
	for all $i,j$. this is called the \textit{frobenius condition}.
\end{thm}
as an example, consider $M=T^{2}$ and $X=2\p_{x}+3\p_{y}$.
representing $T^{2}$ as a quotient of the square, the flow line of $X$ can
be thought of as a line with slope $3/2$.
all of the flow lines is, intuitively, called a \textit{foliation by submflds}.
in particular, the flow lines would ``be" a submanifold.

if on the other hand, we consider $X=\p_{x}+\sqrt{2}\p_{y}$, a single
flow line would densely cover $T^{2}$.
the flow lines also define a foliation, i.e. a locally nice-looking
partition into submflds, but globally they are not submflds.

\begin{lm}[title=Flow-straightening lemma]
	given lin. indep. $X_{1},\dots,X_{r}$ VF on $M$ w $[X_{i},X_{j}]=0$ and
	$p$ s.t. $X_{p}\neq0$, $\ex(U,\vphi)$ near $p$ s.t. $X_{i}=\p_{i}$ wrt the chart.
\end{lm} \

\begin{pf}
	choose a codim $r$ $N$ s.t. $\trm{span}((X_{i})_{p},T_{p}N)=T_{p}M$.
	consider $F:(-\ep,\ep)^{r}\times N\gto M$ given as:
	\begin{equation*}
		F(t_{1},\dots,t_{r},p) \ = \ \Phi_{t_{1}}^{x_{1}}\circ\cdots\circ\Phi_{t_{r}}^{x_{r}}(p)
	\end{equation*}
	note the flows can be rearranged, and we have $TF(\p_{t_{i}})=X_{i}$.
	then $T_{(0,p)}F(s,v)=v+\sum s_{i}(X_{i})_{p}$, so $\Phi$ full rank at $(0,p)$.
	by inv fn thm, we get a chart.
\end{pf}

pf of frebenius thm: forward dir is true since if $X_{i},X_{j}$ are
tangent to submfld $S$, then so is $[X_{i},X_{j}]$.
thus, its in their span.
OTOH, we reduce to $[\cdot,\cdot]=0$ case.
we have (?) in $U\subseteq\bR^{m}$ in coords:
\begin{equation*}
	X_{i}=\sum_{j=1}^{r}a_{ij}(x)\p_{j}+\sum_{j=r+1}^{m}b_{ij}(x)\p_{j}
\end{equation*}
by a lin. change of coords, $a_{ij}(0)=I$ and $b_{ij}(0)=0$.
define $X_{i}=\sum_{j=1}^{r}a_{ij}X'_{j}$.
the $X'_{j}$ span the same thing and commute, so by the lemma,
they give a submfld.

\lecdate{Lec 17 - Mar 12 (Week 9)}
TODO: sort these notes.
... are we reviewing?

to understand VF, need to understand them as:
\begin{itemize}
	\item \textit{coord-wise}: $X=\sum X_{i}\p_{i}$, where $X_{i}$'s are smooth fns
		(in any chart on $M$)
	\item \textit{abstract}: a smooth section of the tangent bundle
		(smooth choice of tangent vec for each point)
	\item \textit{functional} [or sheaf-theoretic?]: a derivation on $C^{\infty}(M)$
	\item \textit{integral}: sth that generates a flow
		(a 1-param family of diffeos)
\end{itemize}

we can do sth similar for diff forms.
\begin{itemize}
	\item \textit{coord-wise}: $k$-form $\omega=\sum f_{I}\di\omega_{i_{1}}\wedge \cdots \wedge \d\omega_{i_{k}}$
	\item \textit{abstract}: $\omega$ is a smooth choice of $k$-covec at each point
		(a section of some [the cotangent] vector bundle)
	\item \textit{functional}: $\omega$ a kind of fn $(\frk{X}(M))^{k}\gto C^{\infty}(M)$
	\item \textit{integral}: a $k$-form is sth integrable over a $k$-mfld
\end{itemize}

\begin{xmp}
	given $f\in C^{\infty}(M)$, defn its (exterior) differential
	$\d f:\frk{X}(M)\gto C^{\infty}(M)$ by:
	\begin{equation*}
		\d f(X) = X(f)
	\end{equation*}
	note, $\d f$ is $C^{\infty}(M)$-linear, i.e.:
	\begin{equation*}
		\d f(gX+hY)=g \d f(X)+h\d f(Y)
	\end{equation*}
	that is an exercise.
\end{xmp} \

\begin{defn}[title=Functional]
	a \textbf{1-form} on $M$ is a $C^{\infty}(M)$-linear map $\frk{X}(M)\gto C^{\infty}(M)$.
	the set ($C^{\infty}(M)$-alg) of 1-forms is denoted $\Omega^{1}(M)$.
\end{defn}
hey, rmbr dual vspaces?
given $\alpha\in V^{*}$, we can write $\alpha(v)$ as $\ang{\alpha,v}$.
this emphasizes that if $V$ finite-dim, $(V^{*})^{*}=V$ canonically.
that is, we can think of $V$ as a hom $V^{*}\gto\bR$.
note $V^{*}\simeq V$, but not canonically unless we first fix an inner prod.

\begin{defn}
	given basis $e_{1},\dots,e_{n}$ of $V$, the \textbf{dual basis} for $V^{*}$
	is $e^{1},\dots,e^{n}$, where $e^{i}(e_{j})=\delta_{ij}$.
\end{defn}
cotangent space is defnd precisely as u think.
its elems are called \textbf{covectors}.
in $T_{p}\bR^{n}$, we have a distinguished basis $\p_{1},\dots,\p_{n}$.
its dual basis is written $\d x_{1},\dots,\d x_{n}$, where:
\begin{equation*}
	\d x_{i}(a_{1}\p_{1}+\dots +a_{n}\p_{n})=a_{i}
\end{equation*}

\begin{defn}
	given a lin. map $F:V\gto W$, theres a \textbf{dual map} $F^{*}:W^{*}\gto V^{*}$
	defn'd by $F^{*}\omega(v)=\omega(Fv)$.
\end{defn}
given smooth map $f:M\gto N$, we have a linear map $T_{p}f:T_{p}M\gto T_{f(p)}N$
given by $v\mto f_{*}v=T_{p}f(v)$.
therefore, $T_{p}^{*}f:T_{p}^{*}M\gto T_{f(p)}^{*}N$ is given as
$v\mto f^{*}v=T_{p}^{*}f(v)$.
the ``base" map is as a pushforward, i.e. vecs are \textit{covariant}.
the dual is as a pullback, i.e. covecs are \textit{contravariant}.

\begin{xmp}
	given $f\in C^{\infty}(M)$, we can defn a covec $\d_{p}f\in T^{*}_{p}M$ by:
	\begin{equation*}
		\d_{p}f(v)=v(f)=D_{v}f
	\end{equation*}
	in particular, in $\bR^{n}$, we can write
	$\d_{p}f=\p_{x_{1}}f\d x_{1}+\dots +\p_{x_{n}}f\d x_{n}$.
	exercise... lol.
\end{xmp} \

\begin{prop}
	given $f\in C^{\infty}(M,N),g\in C^{\infty}(N)$, we have
	$\d_{p}(f^{*}g)=f^{*}\d_{p} g$.
\end{prop} \

\begin{pf}
	for any $v\in T_{p}M$:
	\begin{align*}
		\ang{\d_{p}(f^{*}g),v}
		& \ = \ v(f^{*}g) \\
		& \ = \ v(g\circ f) \\
		& \ = \ (T_{p}f(v))(g) \\
		& \ = \ \ang{\d_{f(p)}g,T_{p}f(v)} \\
		& \ = \ \ang{f^{*}\d_{f(p)}g,v}
	\end{align*}
	note $\ang{\cdot,\cdot}$ is not an inner prod here.
\end{pf}

\lecdate{Lec 18 - Mar 17 (Week 10)}
last time, we talked abt covecs - this is enough for a ``vibe" of a defn.
what is a 1-form?
it is a smooth choice of covec at each pt of a mfld.

\begin{defn}[title=Coordinate-wise]
	a map $p\mto\alpha_{p}$ sending pts to covecs is a \textbf{diff. 1-form} if
	for every chart $(U,\vphi)$:
	\begin{equation*}
		\alpha_{\vphi(p)} \ = \ \sum_{i=1}^{m}g_{i}(\vphi(p))\di x_{i}
	\end{equation*}
	where $g_{i}:\vphi(U)\gto\bR$ is smooth.
\end{defn}

for instance:
\begin{enumerate}[1)]
	\item $f:M\gto\bR$ as $\d f = \sum\p_{x_{i}}f-\d x_{i}$ on a chart.
		as an exercise, show this is indep of choice of chart.
	\item $y\d x$ on $\bR^{2}$ is not $\d f$ for any $f$.
		to see this, note its partials don't commute [why?], or that
		$y\d x$ depends on $y$ so $f$ depends on $y$, so $\p_{y}f\neq0$.
		[what?]
\end{enumerate}

\begin{defn}[title=Functional]
	a 1-form is a $C^{\infty}(M)$-linear $\frk{X}(M)\gto C^{\infty}(M)$
	given by applying the covec to pts.
\end{defn} \

\begin{prop}
	let $\alpha$ be a 1-form in the functional sense.
	then for any VF $X$, there exists a well-defn'd covec $\alpha_{p}$ w $p\in M$ s.t.
	$\alpha(X)(p)=\alpha_{p}(X_{p})$.
\end{prop} \

\begin{pf}
	wts if $X_{p}=Y_{p}$, then $\alpha(X)(p)=\alpha(Y)(p)$.
	therefore, this gives a well-defn'd map $X_{p}\mto\alpha_{p}(X_{p})$.
	that it's linear, ie a covec, is an exercise.
	by linearity, enough to show that if $X_{p}=0$ then $\alpha(X)(p)=0$.

	indeed, if $X_{p}=0$, write $X=\sum f_{i}Y_{i}$ where $f_{i}(p)=0$
	($Y_{i}=\p_{i}$ on a chart near $p$, then extend somehow).
	thus:
	\begin{equation*}
		\alpha(X)(p) \ = \ \sum_{i}f_{i}(p)\alpha(Y_{i})(p) \ = \ 0
	\end{equation*}
\end{pf} \

\begin{prop}\vspace{-0.2in}
	\begin{enumerate}[a)]
		\item functional 1-forms can be restricted to a form on an open
			set $U\subseteq M$.
		\item the two defns are equiv.
	\end{enumerate}
\end{prop}
note that this is non-trivial since a VF on $U$ is not always a restriction of
a VF on $M$.
\begin{exr}
	state carefully what the above means, then prove it.
\end{exr}

given $g:N\gto\bR$, recall a pullback of $g$ along $f:M\gto N$:
\begin{equation*}
	f^{*}g \ = \ g\circ f
\end{equation*}
for covecs, there is an induced map $T_{p}f:T_{p}M\gto T_{f(p)}N$ by
$v\mto f_{*}v$, which in turn induces a map $T_{p}^{*}N\gto T_{p}^{*}M$ by
$\nu\mto f^{*}\nu$, where $\nu$ is a covec.

now, given a 1-form $\omega\in\Omega^{1}(N)$, we can define its \textbf{pullback} along $f$ by:
\begin{equation*}
	(f^{*}\omega)_{p} \ = \ f^{*}(\omega_{f(p)})
\end{equation*}
we claim this is indeed a 1-form, i.e. its a smooth choice.
indeed, in coords, we have a smooth fn $f:U\gto V$ where
$U,V\subseteq\bR^{m},\bR^{n}$ resp.
if $\omega\in\Omega^{1}(V)$ given by $\sum g_{i}\d x_{i}$, then $T_{p}f$ has smooth
coord fns $\p_{x_{j}}f_{i}$. note:
\begin{equation*}
	f^{*}\big\rvert_{f(p)} \ = \ (T_{p}f)^{\intercal}
\end{equation*}
so:
\begin{equation*}
	(f^{*}\omega)_{p} \ = \ f^{*}\bigg\rvert_{f(p)}(\omega_{f(p)})
	\ = \ f^{*}\bigg\rvert_{f(p)}\sum g_{i}(f(p))\di x_{i}
\end{equation*}
these are all smooth in $p$, so we are done.

now, some properties of pullbacks:
\begin{itemize}
	\item $(fg)^{*} \ = \ g^{*}f^{*}$
	\item if $g:M\gto\bR$ smooth, $f^{*}(g\omega) \ = \ f^{*}g \cdot f^{*}\omega$
	\item $f^{*}(\d g) \ = \ \d(f^{*}g)$
	\item in coords:
		\begin{align*}
			(f^{*}\omega)_{p} \ = \ f^{*}\bigg\rvert_{f(p)}\sum_{i}g_{i}(f(p))\di x_{i}
			& \ = \ \sum_{i}g_{i}(f(p))f^{*}\bigg\rvert_{f(p)}\di x_{i} \\
			& \ = \ \sum_{i}\sum_{j}g_{i}(f(p))\frac{\p f_{i}}{\p x_{j}}\di x_{j}
		\end{align*}
\end{itemize}

\begin{defn}
	the \textbf{cotangent bundle} of a mfld $M$ is a mfld structure on:
	\begin{equation*}
		T^{*}M \ := \ \bigsqcup_{p\in M}T_{p}^{*}M
	\end{equation*}
	with the canonical proj $\pi:T^{*}M\gto M$.
	given $(U,\vphi)$, we get a chart on $T^{*}M$:
	\begin{equation*}
		(\pi^{-1}(U),(\vphi^{-1})^{*}:\pi^{-1}U\gto\vphi(U)\times\bR^{m})
	\end{equation*}
	here, $\bR^{m}$ is a copy of the cotangent space.
	note transition maps are smooth, as they are given by:
	\begin{equation*}
		(\vphi\circ\psi^{-1})^{*} = (T_{p}(\vphi\circ\psi^{-1}))^{\intercal} \ : \ \vphi(U\cap V)\times\bR^{m}\gto\psi(U\cap V)\times\bR^{m}
	\end{equation*}
\end{defn}
this is just repackaging the coord-wise defn.

NOTE: with tangent bundles, $f:M\gto N$ defined $Tf:TM\gto TN$.
here, we can only do sth similar when $f$ is a diffeo, as pts go forward but
vecs go back.
OTOH, if $\omega\in\Omega^{1}(N)$ is defn'd as $s_{\omega}:N\gto T^{*}N$, then
$f^{*}\omega$ is given by:
\begin{equation*}
	s_{f^{*}\omega}(p) \ = \ f^{*}\big\rvert_{f(p)}(s_{\omega}\circ f(p))
\end{equation*}

\begin{defn}[title=Abstract]
	a 1-form is a section of the cotangent bundle
\end{defn}

\subsection{Integration}
in desperate need of more headers rn

given $\gamma:I\gto M$, $\alpha\in\Omega^{1}(M)$, take $\gamma^{*}\alpha\in\Omega^{1}(I)$ as
$\gamma^{*}(\alpha)=f(t)\di t$ for some $f(t)$.

\begin{defn}
	we define:
	\begin{equation*}
		\int_{\gamma}\alpha \ = \ \int_{a}^{b}f(t)\di t
	\end{equation*}
	where $I=[a,b]$.
	we'll take smoothness on closed intervals in the extension way.
\end{defn} \

\begin{prop}
	if $\alpha=\d f$ for some $f:M\gto\bR$, then:
	\begin{equation*}
		\int_{\gamma}\alpha \ = \ f(\gamma(b))-f(\gamma(a))
	\end{equation*}
\end{prop} \

\begin{pf}
	we have that
	$\gamma^{*}\d f \ = \ \d(\gamma^{*}f) \ = \ \d(f\circ\gamma) \ = \ \p_{t}(f\circ\gamma)\d t$.
	apply usual ftc.
\end{pf}

for instance, $y\d x$ is not $\d f$, i.e. an \textit{exact form}, since
by integrating counterclockwise on $\p[0,1]^{2}$, we get:
\begin{equation*}
	\int_{\gamma}y\di x \ = \ 0+0+\int_{0}^{1}1(-\d t) + 0 \ = \ 1 \ \neq \ 0 \ = \ f(0,0)-f(0,0)
\end{equation*}

\begin{prop}
	integrating is invariant wrt reparametrization:
	\begin{equation*}
		\int_{\gamma}\alpha \ = \ \pm\int_{\gamma\circ\kappa}\alpha
	\end{equation*}
	where $\kappa:[c,d]\gto[a,b]$ is a diffeo, and $\pm$ depends on orientation.
\end{prop} \

\begin{pf}
	we have:
	\begin{equation*}
		\int_{\gamma\circ\kappa}\alpha \ = \ \int_{c}^{d}(\gamma\circ\kappa)^{*}\alpha
		\ = \ \int_{c}^{d}\kappa^{*}\gamma^{*}\alpha \ = \ \pm \int_{\kappa(c)}^{\kappa(d)}\gamma^{*}\alpha
	\end{equation*}
	after a substitution.
\end{pf}

recall $\alpha$ is \textit{exact} if its $\d f$ for some $f\in C^{\infty}$.
in this case, $\int_{\gamma}\alpha$ depends only on endpts.

\textit{fact: this is an iff}.
if $\int_{\gamma}\alpha$ depends only on endpts, then defn $f(p)=\int_{\gamma}\alpha$ where
$\gamma$ is any curve $p_{0}$ to $p$ (where we can choose $p_{0}$).
to ensure $f$ smooth and $\d f=\alpha$, do a coord calculation.
ofc, this is iff $\int_{\gamma}\alpha=0$ for $\gamma$ a loop.

\begin{xmp}
	let $M=S^{1},\alpha=\d\theta$. then:
	\begin{equation*}
		\int_{S^{1}}\d\theta \ = \ \int_{\gamma}\d\theta \ = \ 2\pi
	\end{equation*}
	where $\gamma$ inj on $(a,b)$ and $\gamma(a)=\gamma(b)$.
	thus, $\d\theta$ is not an exact form.
\end{xmp}
